\section{Decodificador 3x8}

\subsection{Tabela Verdade e Expressões Booleanas}

A tabela verdade de um decodificador 3x8 com entrada em binário puro ($A_2, A_1, A_0$) e entrada de habilitação ($En$) é apresentada a seguir. Quando $En = 0$, o decodificador está desabilitado e todas as saídas permanecem em nível baixo (0), independentemente dos valores de entrada. Quando $En = 1$, a combinação de entrada ativa exclusivamente a saída correspondente.

\begin{table}[H]
    \centering
    \caption{Tabela Verdade do Decodificador 3x8 com Enable}
    \begin{tabular}{cccc|cccccccc}
        \toprule
        \textbf{$En$} & \textbf{$A_2$} & \textbf{$A_1$} & \textbf{$A_0$} & \textbf{$Y_0$} & \textbf{$Y_1$} & \textbf{$Y_2$} & \textbf{$Y_3$} & \textbf{$Y_4$} & \textbf{$Y_5$} & \textbf{$Y_6$} & \textbf{$Y_7$} \\
        \midrule
        0             & X              & X              & X              & 0              & 0              & 0              & 0              & 0              & 0              & 0              & 0              \\
        1             & 0              & 0              & 0              & 1              & 0              & 0              & 0              & 0              & 0              & 0              & 0              \\
        1             & 0              & 0              & 1              & 0              & 1              & 0              & 0              & 0              & 0              & 0              & 0              \\
        1             & 0              & 1              & 0              & 0              & 0              & 1              & 0              & 0              & 0              & 0              & 0              \\
        1             & 0              & 1              & 1              & 0              & 0              & 0              & 1              & 0              & 0              & 0              & 0              \\
        1             & 1              & 0              & 0              & 0              & 0              & 0              & 0              & 1              & 0              & 0              & 0              \\
        1             & 1              & 0              & 1              & 0              & 0              & 0              & 0              & 0              & 1              & 0              & 0              \\
        1             & 1              & 1              & 0              & 0              & 0              & 0              & 0              & 0              & 0              & 1              & 0              \\
        1             & 1              & 1              & 1              & 0              & 0              & 0              & 0              & 0              & 0              & 0              & 1              \\
        \bottomrule
    \end{tabular}
\end{table}

A partir da tabela, extraímos as expressões booleanas simplificadas para cada saída. Como o decodificador só funciona quando habilitado, o termo $En$ deve estar presente (em nível lógico alto) em todas as equações, multiplicando os mintermos das entradas:

\begin{align*}
    Y_0 & = En \cdot \overline{A_2} \cdot \overline{A_1} \cdot \overline{A_0} \\
    Y_1 & = En \cdot \overline{A_2} \cdot \overline{A_1} \cdot A_0            \\
    Y_2 & = En \cdot \overline{A_2} \cdot A_1 \cdot \overline{A_0}            \\
    Y_3 & = En \cdot \overline{A_2} \cdot A_1 \cdot A_0                       \\
    Y_4 & = En \cdot A_2 \cdot \overline{A_1} \cdot \overline{A_0}            \\
    Y_5 & = En \cdot A_2 \cdot \overline{A_1} \cdot A_0                       \\
    Y_6 & = En \cdot A_2 \cdot A_1 \cdot \overline{A_0}                       \\
    Y_7 & = En \cdot A_2 \cdot A_1 \cdot A_0
\end{align*}